\section{Limitations}
\label{sec:limitations}

All reported values are best fits; no posterior intervals are claimed. Endpoint posterior uncertainty is not propagated into the difference fields. C1 and C2 are not statistically independent constructions: both use Planck information. R carries $r_{d}=135.39$~Mpc while the C1/C2 references use recomputed CAMB rulers near $147.1$~Mpc (Table~\ref{tab:calibration}). The constrained reconstruction domain is $z\le1.8$; the order-3 continuation fails at the $z=2.33$ Ly$\alpha$ anchor. The calibrator result covers the declared contract only ($\Delta\mu_{\mathrm{cal}} = 0$, a qualitative assumption, not quantitatively established).

\begin{table}[!htbp]
  \centering
  \caption{Limitations and current status.}
  \label{tab:limitations}
  \begin{tabular}{ll}
    \toprule
    Issue & Current status \\
    \midrule
    Endpoint posterior uncertainty & Not propagated \\
    C1/C2 statistical independence & Shared Planck information \\
    Order-selection bootstrap & $P = 0.216$ \\
    Constrained reconstruction range & $z \le 1.8$ \\
    Ly$\alpha$ continuation test & Fails at $z = 2.33$ \\
    \bottomrule
  \end{tabular}
\end{table}

The distance-sector endpoint and the Planck-compatible
references prefer quantitatively different expansion histories
($H_{0}$ \DistanceHZero{} vs \CmbOneHZero{}, with endpoint-specific
acoustic rulers per Table~\ref{tab:calibration}).

\subsection{Relation to neighbouring approaches}
\label{sec:literature}

DESI DR2 BAO remain compatible with flat $\Lambda$CDM in isolation,
while combinations with CMB and supernovae can favor evolving dark
energy of the $w_{0}w_{a}$ form~\citep{DESIDR2,
Chevallier2001,Linder2003}. Our reconstruction is not a dark-energy equation of
state: it is a reconstruction of the
required observational transformation between two calibrated
histories, with dimensionality selected through held-out data rather
than by allowing enough function freedom to improve the joint fit.
That selection discipline is also what separates this work from
unconstrained Gaussian-process or smoothing reconstructions of the
expansion history~\cite{Shafieloo2006} and from generic
cosmography~\cite{Cattoen2007}, which do not penalize flexibility by
prediction. Early-dark-energy solutions modify the pre-recombination
universe and its perturbations; the reconstruction assumes no recombination
change and carries no perturbation sector. For the broader solution
landscape we refer to the community review~\cite{DiValentino2021}.
This work does not introduce another
equation-of-state fit. It compares separately calibrated histories
and asks whether their difference is compact and
reference-insensitive within the two tested Planck-compatible
constructions. This decomposition is complementary to the scale-vs-shape
obstruction argument \citep{Zhou2026}: uniform sound-horizon rescaling
changes scale but cannot remove the BAO--SN shape disagreement, and our
R endpoint contains both a large ruler shift ($\sim8\%$) and a
redshift-dependent shape difference.
